% Chapitre 2 : écran fictif d'une issue avec sa liste de cases.
\begin{figure}[tbp]
\centering
\begin{tikzpicture}[x=1cm,y=1cm,
  texte/.style={font=\sffamily\scriptsize, anchor=west, inner sep=0pt}]
  \fenetre{12.4}{7.2}{Issue \#89}
  \node[font=\sffamily\bfseries, anchor=west] at (0.35,6.15)
    {Exposé I : les énoncés précis et leur littérature \mdseries\color{black!50}\#89};
  \node[etiquette, fill=tonregle, text=white, anchor=west] at (0.4,5.55) {Ouverte};
  \node[font=\sffamily\scriptsize, text=black!65, anchor=west] at (1.55,5.55)
    {2 tâches faites sur 4};
  \fill[black!10] (4.4,5.5) rectangle (6.4,5.6);
  \fill[tonregle] (4.4,5.5) rectangle (5.4,5.6);
  \draw[black!20] (0.35,5.2) -- (12.05,5.2);
  % corps de l'issue
  \node[texte, anchor=north west, text width=7.4cm, align=left] at (0.4,4.95)
    {Relever les énoncés de l'exposé I et chercher chacun dans la littérature, avec la
    page. Suite de \textcolor{tonregle}{\#84}.};
  \foreach \y/\f/\t in {3.95/1/lire l'exposé I{,} chercher les énoncés précis,
      3.45/1/recherche bibliographique{,} énoncé par énoncé,
      2.95/0/formaliser ce qui n'est pas dans la littérature,
      2.45/0/proposer le résultat à la bibliothèque commune} {
    \draw[black!45, rounded corners=1pt] (0.45,\y-0.13) rectangle (0.71,\y+0.13);
    \ifnum\f=1
      \fill[tonregle, rounded corners=1pt] (0.45,\y-0.13) rectangle (0.71,\y+0.13);
      \draw[white, line width=0.9pt, line cap=round, line join=round]
        (0.5,\y) -- (0.57,\y-0.07) -- (0.67,\y+0.08);
      \node[texte, text=black!55] at (0.9,\y) {\t};
    \else
      \node[texte] at (0.9,\y) {\t};
    \fi
  }
  % fil de l'issue
  \draw[black!20] (0.35,1.95) -- (8.0,1.95);
  \fill[black!35] (0.58,1.45) circle (0.07);
  \node[texte, text=black!65] at (0.85,1.45)
    {référencée par le commit \texttt{\color{tonregle}4233568}, il y a deux jours};
  \fill[black!35] (0.58,0.85) circle (0.07);
  \node[texte, text=black!65] at (0.85,0.85)
    {étiquette \emph{en attente} ajoutée : la page 12 attend une réponse de l'archive};
  % colonne de droite
  \node[font=\sffamily\scriptsize\bfseries, text=black!65, anchor=west] at (8.6,4.8) {Étiquettes};
  \node[etiquette, draw=black!40, fill=black!5, anchor=west] at (8.6,4.35) {lecture};
  \node[etiquette, draw=black!40, fill=black!5, anchor=west] at (9.6,4.35) {Lean};
  \node[etiquette, draw=tonregle, fill=tonregle!8, text=tonregle, anchor=west] at (10.45,4.35)
    {en attente};
  \node[font=\sffamily\scriptsize\bfseries, text=black!65, anchor=west] at (8.6,3.7) {Liens};
  \node[texte, text=tonregle] at (8.6,3.3) {\#84 exposé I, notations};
  \node[texte, text=tonregle] at (8.6,2.9) {\#95 rédactions en Lean};
  \node[font=\sffamily\scriptsize\bfseries, text=black!65, anchor=west] at (8.6,2.25) {Projet};
  \node[texte] at (8.6,1.85) {Lecture des exposés};
\end{tikzpicture}
\caption{Écran fictif : une issue. Le titre pose la question, la liste de cases en montre
l'avancement, et le fil garde la trace du commit qui l'a fait avancer et de ce qu'on
attend.}
\label{fig:ecran-issue}
\end{figure}
